\section{Integer multiplication on fixed tapes}\label{sec:integers}

Write each input as a polynomial evaluated at a power of two. To
multiply the integers it suffices to convolve the coefficient lists
exactly and propagate carries. We compute this convolution with two
forward Fourier transforms, one coordinatewise product and one
inverse transform. On tapes, the choice of the transform length
determines how much time is spent reaching the records that each
operation needs.

\begin{theorem}[Multiplication at a fixed exponent]\label{thm:integer}
For each \(\epsilon>0\) there is a deterministic multitape Turing
machine, with fixed alphabet and tape count, that computes the full
product of two \(n\)-bit integers in worst-case time
\(O_\epsilon(n(1+\log n)^{1/2+\epsilon})\), visiting
\(O_\epsilon(n)\) cells on its tapes.
\end{theorem}

Fix \(\epsilon>0\), and choose a rational exponent
\(\tfrac12<\rho<\min\{1,\tfrac12+\epsilon\}\) for
Theorem~\ref{thm:interchange} and Lemma~\ref{lem:tape-walsh}.
We will store the digits in \(N=2^k\) records of width \(O_\rho(k)\),
with \(k=\Theta(\log n)\) and \(Nk=\Theta(n)\). We arrange the
convolution on many short cyclic axes. Each transform then needs
only \(O(\log k)\) layers of record permutations, phase
multiplications and selected-bit Walsh transforms, preserving
the Euclidean norm at every layer boundary. The time budget for
Theorem~\ref{thm:integer} is
\begin{equation}\label{eq:integer-budget}
 O_\rho\!\left(Nk(1+k)^\rho(1+\log(2+k))^{O(1)}\right).
\end{equation}
Products within a single record of \(O(k)\) bits take
\(O(k\log(2+k))\) steps with the multitape multiplier of
Harvey and van der Hoeven \cite{hvh}. Marking the extents of its
scratch areas lets us clear them and return the heads within the
same time bound.

\subsection{Packing the digits into short cyclic axes}

For \(n=1\) use direct multiplication. For \(n\ge2\), let
\(X(t)=\sum_{u<h}x_u t^u\) and \(Y(t)=\sum_{u<h}y_u t^u\)
be the radix-\(2^b\) expansions of the inputs, with
\(0\le x_u,y_u<2^b\). Set
\[
 b=\lceil\log_2 n\rceil,\qquad h=\lceil n/b\rceil,\qquad
 k=\lceil\log_2(4h)\rceil,\qquad N=2^k.
\]
A cyclic convolution of length at least \(2h-1\) recovers every
coefficient of \(X(t)Y(t)\). We choose its length \(S\) as a product
of pairwise coprime factors \(a_i\). Factor \(a_i\) has an address
field of \(l_i\) bits, so its padded axis has \(2^{l_i}\) positions.
We require
\begin{equation}\label{eq:integer-box}
 \sum_i l_i=k,\qquad \max_i l_i=O(\log k),\qquad
 d=O(k/\log k),\qquad S:=\prod_{i=1}^d a_i\ge 2h .
\end{equation}
Short axes give polynomial-size tables in \(k\). Filling a constant
fraction of the padded box ensures that these tables can act on a
stream of only \(O(h)\) records, even with separate binary address
fields for all \(d\) axes.

For all sufficiently large \(k\), put \(l=20\lceil\log_2 k\rceil\)
and \(x=2^l\). Take \(\lfloor k/l\rfloor\) distinct odd primes
from \([x(1-k^{-2}),x]\), each with \(l_i=l\). If \(l\) does not
divide \(k\), append \(a_d=2^{k\bmod l}\), with \(l_d=k\bmod l\).
Baker, Harman and Pintz \cite{bhp} give a prime in
\([y-y^{0.6},y]\) for all sufficiently large \(y\). The chosen
interval contains \(k\) disjoint intervals of this form with right
endpoints near \(x\), since \(kx^{0.6}=o(x/k^2)\). Enumeration
and trial division take \(\operatorname{poly}(k)\) bit steps because
\(x=k^{O(1)}\). The resulting volume satisfies
\[
 S\ge N(1-k^{-2})^{\lfloor k/l\rfloor}\ge N/2\ge 2h.
\]
The machine tests the size and coprimality conditions and uses
direct multiplication on the bounded initial range where they fail.

Store the \(N\) positions in lexicographic order on the padded
axes \(d,d-1,\ldots,1\), with axis \(1\) varying fastest. A valid
address has digits \(t_i<a_i\) for every \(i\). In tape order, these
valid addresses enumerate the ranks \(0,\ldots,S-1\) in mixed radix.
Count the input length with an amortized binary counter and reverse
the streams to read their low bits first. As the machine enumerates
the box it places the successive radix-\(2^b\) digits at the first
\(h\) valid addresses, writing zero at every other position. A
simultaneous scan of the \(k\) address bits and the marked fields
containing their limits tests validity in \(O(k)\) steps. The writes,
scans and rewinds together cost \(O(Nk)\).

\subsection{Changing from digit order to residue order}

The digits enter the tape in increasing order of their mixed-radix
rank. For convolution the indices must add coordinatewise modulo the
\(a_i\). Put \(P_i=\prod_{j<i}a_j\) and
\(\mu_i=P_i^{-1}\bmod a_i\). If \(t=\sum_iP_it_i\) has digits
\(t_i\), use the scaled residue
\begin{equation}\label{eq:integer-crt}
 y_i=\mu_i t\bmod a_i
     =t_i+\mu_i\sum_{j<i}P_jt_j\bmod a_i .
\end{equation}
These coordinates give an isomorphism
\(\mathbb Z/S\mathbb Z\to\prod_i\mathbb Z/a_i\mathbb Z\).
The scaling by \(\mu_i\) makes the coefficient of the current digit
one. The permutation can therefore use rotations controlled by
earlier digits, as in the prime-factor algorithm \cite{good} and the
triangular layout of \cite[Equation~(42)]{Tape}.

Build a balanced binary tree of consecutive axis intervals, splitting
an interval after the first \(\lfloor |I|/2\rfloor\) axes. The map for
\(I\) is Equation~\eqref{eq:integer-crt} with the sum restricted to
\(j\in I\). A leaf is the identity. Once the two child maps have been
applied, each right-child axis \(i\) only needs the rotation
\[
       \mu_i\sum_{j\ \mathrm{in\ the\ left\ child}} P_jt_j
       \pmod {a_i}.
\]
The current left-child address determines its original digits \(t_j\)
by solving the triangular equations in increasing order. Use zero
rotation for an invalid left-child address and fix values outside
\([0,a_i)\) on each target axis. Each step is then a permutation of
the full padded box. For the inverse map, reverse the depth order
and negate the rotations.

Lemma~\ref{lem:controlled-blocks} combines all disjoint intervals at
one depth as a single controlled permutation. The control
bits lie in the left children, which collectively contain at most
half the axes. Since all fields except possibly the last have equal
width, the controls occupy at most \(k/2+O(\log k)\) bits. The target
blocks have width \(O(\log k)\). For a given control setting we can
recover the left digits and generate all local rotation tables in
\(\operatorname{poly}(k)\) steps. The integers \(P_j\) have at most
\(k+1\) bits and a table has at most \(2^l=\operatorname{poly}(k)\)
entries. The lemma charges for generating these tables as well as for
moving the records. A constant fraction of address bits lies outside
the controls, so the lemma applies at every depth. The \(O(\log k)\)
depths fit the budget in Equation~\eqref{eq:integer-budget} in both
directions.

\subsection{Short Fourier transforms on occupied fibers}

After the CRT permutation, each factor has length \(a_i\), but its
address occupies \(l_i\) bits. A radix-two chirp computation may need
several times as many coordinates as fit along that axis. Padding
every axis by such a factor would increase the size of the tensor
product exponentially in the number of axes. Instead we borrow
coordinates from adjacent fibers, perform a short transform, and
restore the borrowed values exactly in the ideal complex arithmetic.

Write \(e(t)=\exp(2\pi i t)\), and put
\[
 U_a=\bigl(a^{-1/2}e(jh/a)\bigr)_{0\le j,h<a},\qquad
 H=2^{-1/2}\begin{pmatrix}1&1\\1&-1\end{pmatrix}.
\]
\begin{lemma}[Synchronized short Fourier transforms]\label{lem:synchronized-short}
For the axis lengths and widths of Equation~\eqref{eq:integer-box},
let \(k\) be sufficiently large. On the \(N=2^k\) coordinates of the
padded residue array, both
\[
             \mathcal U=\bigotimes_i(U_{a_i}\oplus I_{2^{l_i}-a_i})
             \qquad\text{and}\qquad \mathcal U^*
\]
factor into \(O(\log(2+k))\) unitary layers. Partition the axes into
disjoint consecutive blocks of \(O(\log(2+k))\) address bits each.
Every layer has one of the following forms.
\begin{enumerate}
\item A permutation acting independently within each block.
\item A diagonal whose value at an address is the product of one
      unit-modulus coefficient from each block.
\item A normalized Walsh transform on one designated bit from each
      participating block, using \(O(k/\log(2+k))\) bits in total.
\end{enumerate}
For any fixed \(C>0\), a deterministic fixed-tape machine constructs
the schedule, the block permutations, and the block phase tables
at \(Ck+O(\log(2+k))\) fractional bits in \(\operatorname{poly}_C(k)\)
steps. It can then stream the coefficients of any full diagonal to
absolute error \(2^{-Ck}\) in \(O_C(Nk\log(2+k))\) steps.
\end{lemma}

\begin{proof}
It suffices to construct the forward transform. Reversing the schedule,
inverting the permutations and conjugating the phases gives its adjoint.

\paragraph{An exact transform on occupied coordinates.}
Fix an axis of length \(a\), and take the least power of two
\(r\ge 2a-1\). Inside \(\mathbb C^{2r}\) let \(T\) be the space of
the first \(a\) coordinates and \(P\) its orthogonal projection.
We construct \(U_a\oplus I_{2r-a}\) in \(O(\log a)\) layers of
permutations, phase diagonals and operations on disjoint pairs.
The inequality \(2r<8a\) will bound the number of neighboring values
to borrow.

A truncated chirp gives a unitary \(V\) whose compression to \(T\)
is a fixed positive multiple of \(U_a\). In detail, zero-extend
\(a^{-1/2}e(-h^2/(2a))\), \(|h|<a\), to a cyclic kernel of length
\(r\). Its Fourier multipliers are
\[
              \lambda_\xi
        =a^{-1/2}\sum_{|h|<a} e(-h^2/(2a)+\xi h/r).
\]
We show below that \(|\lambda_\xi|\le40\), uniformly in \(a,\xi\).
Choose \(J=64\) and \(\phi=\pi/(4J+2)\). With
\(z=\lambda_\xi\sin\phi\) we have \(|z|\le40\pi/258<1/2\).
Apply \(U_r\) on the first \(r\) coordinates, act on each pair
\((\xi,r+\xi)\) by the unitary
\[
             L(z)=
             \begin{pmatrix}z&\sqrt{1-|z|^2}\\
                            \sqrt{1-|z|^2}&-\overline z\end{pmatrix},
\]
then apply \(U_r^*\) on the first \(r\) coordinates. The compression
to these \(r\) coordinates is convolution by the truncated chirp
times \(\sin\phi\). The differences \(j-h\) of two indices of \(T\)
lie in \((-a,a)\), so their representatives in this interval are
unambiguous modulo \(r\). Multiplying the
active coordinates by \(e(j^2/(2a))\) before and after the
convolution therefore gives Bluestein's factorization
\cite{bluestein}. For the resulting unitary \(V\) we have
\(PV|_T=\sin\phi\,U_a\).

Amplitude amplification now makes the first block exact
\cite{BHMT,berry,MZ}. For each vector \(v\) of an orthonormal basis
of \(T\), put \(w_v=(I-P)Vv/\cos\phi\). These vectors in \(T^\perp\)
are orthonormal because both \(V\) and \(U_a\) are isometries.
In the plane spanned by \(w_v,U_av\),
\[
                    Vv=\cos\phi\,w_v+\sin\phi\,U_av .
\]
The reflection \(I-2P\) fixes \(w_v\) and reverses \(U_av\). On this
plane \(2VPV^*-I\) reflects across \(Vv\), since the other vectors
of the basis of \(V(T)\) are orthogonal to the plane. The product
of the reflections rotates through \(2\phi\) towards \(U_av\).
Consequently
\[
                  B=\bigl[(2VPV^*-I)(I-2P)\bigr]^J V
                    = U_a\oplus B'
\]
for a unitary \(B'\) on \(T^\perp\), since \((2J+1)\phi=\pi/2\).
The only reflection we must implement in the original coordinates
is a diagonal sign change.

To restore the neighboring values, we use \(B\) as a change of basis.
The chirp
\[
             \chi(t)=e\bigl((t^2+(a\bmod 2)t)/(2a)\bigr)
\]
is \(a\)-periodic, and \(\chi(j)\chi(-h)/\chi(j-h)=e(jh/a)\).
Thus \(U_a=D_1 C D_3\), where \(D_1(j)=\chi(j)\), \(D_3(h)=\chi(-h)\)
and \(C\) is the circulant with kernel \(a^{-1/2}\overline{\chi(t)}\).
In the Fourier basis \(D_2=U_aCU_a^*\) is a diagonal of phases,
\[
               D_2(\xi)=a^{-1/2}
                    \sum_{h=0}^{a-1}\overline{\chi(h)}e(\xi h/a).
\]
Extending each of the three diagonals by the identity gives
\[
 (D_1\oplus I)B^*(D_2\oplus I)B(D_3\oplus I)=U_a\oplus I.
\]
Indeed the middle conjugation acts by \(C\) on \(T\) and by
\((B')^*B'=I\) on its complement. The identity holds on arbitrary
data in the borrowed coordinates (Figure~\ref{fig:borrow}).
Each use of \(U_r\) or its adjoint has radix-two depth \(O(\log a)\)
\cite{cooleytukey}, and the number of such uses is absolute.

\begin{figure}[!htbp]
\centering
\begin{tikzpicture}[
   font=\small, >=Stealth,
   amp/.style={line width=.7pt,-{Stealth[length=1.7mm]}},
   layer/.style={minimum width=9mm,minimum height=6.5mm,
                 draw=ink,inner sep=1.5pt,fill=white,font=\small},
   result/.style={layer,fill=sea!9,draw=sea}
]
\begin{scope}
 \draw[ink!45] (2.4,0)--(0,0)--(0,2.2);
 \draw[amp,ink!65] (0,0)--(2.03,0)
       node[below right,inner sep=1pt] {$w_v$};
 \draw[amp,sea] (0,0)--(0,2.03) node[above,inner sep=2pt] {$U_av$};
 \draw[amp,rust] (0,0)--(18:2.03) node[right,inner sep=2pt] {$Vv$};
 \draw[ink!55,dashed] (0,0)--(54:2.03);
 \draw[sea,amp] (18:1.7) arc[start angle=18,end angle=54,radius=1.7];
 \draw[sea,amp] (54:1.7) arc[start angle=54,end angle=90,radius=1.7];
 \node[font=\footnotesize,fill=white,inner sep=.5pt] at (9:1.21) {$\phi$};
 \node[font=\footnotesize,inner sep=.5pt] at (37:1.25) {$2\phi$};
 \node[font=\footnotesize,inner sep=.5pt] at (73:1.25) {$2\phi$};
\end{scope}
\begin{scope}[shift={(4.05cm,0cm)}]
 \foreach \y/\name in {1.55/T,.65/{T^\perp}}{
    \node[left,inner sep=2.5pt] at (0,\y) {$\name$};
    \draw[ink!65,-{Stealth[length=1.6mm]}] (0,\y)--(6.95,\y);
 }
 \foreach \x/\fiberop/\borrowop in {
     .55/{D_3}/I,1.8/{U_a}/{B'},3.05/{D_2}/I,4.3/{U_a^*}/{(B')^*},5.55/{D_1}/I}{
    \node[layer,fill=sea!4] at (\x,1.55) {$\fiberop$};
    \node[layer] at (\x,.65) {$\borrowop$};
 }
 \node[text=rust] at (1.8,2.15) {$B$};
 \node[text=rust] at (4.3,2.15) {$B^*$};
 \node[result] at (7.55,1.55) {$U_a$};
 \node[result] at (7.55,.65) {$I$};
 \node[font=\footnotesize,text=ink] at (3.05,.04) {Execution order};
\end{scope}
\end{tikzpicture}
\caption{Building \(U_a\oplus I\) on occupied coordinates. In each
plane spanned by \(w_v\) and \(U_av\), a pair of reflections rotates
\(Vv\) towards \(U_av\) by \(2\phi\) (shown with \(J=2\)). On the
right, the cyclic chirp multiplies by \(D_2\) in the amplified basis.
The two transformations of the borrowed coordinates cancel.}
\label{fig:borrow}
\end{figure}


\paragraph{The uniform multiplier bound.}
Set \(u_h=e(-h^2/(2a)+\xi h/r)\) and
\(\alpha_h=\xi/r-(h+1/2)/a\), so that \(u_{h+1}=u_he(\alpha_h)\).
Over \(|h|<a\) the \(\alpha_h\) traverse an interval of length
less than two in steps of \(1/a\). With \(\delta=1/(4\sqrt a)\),
at most three \(\delta\)-neighborhoods of integers meet that interval.
They contain at most \(3(2a\delta+2)\) terms, each of modulus one.
The remaining terms split into at most four consecutive runs with
\(\alpha_h\in[j+\delta,j+1-\delta]\) for some integer \(j\).
On each run write \(u_h=(u_h-u_{h+1})w_h\), where
\[
             w_h=(1-e(\alpha_h))^{-1}
                =\tfrac12+\tfrac i2\cot(\pi\alpha_h).
\]
Here \(|w_h|\le(2\sin(\pi\delta))^{-1}\) and the imaginary parts
are monotone. Summation by parts bounds the sum on the run by the
endpoint moduli plus the variation of \(w_h\), at most
\(2/\sin(\pi\delta)\le1/\delta\). Adding the contributions and
normalizing gives
\[
            |\lambda_\xi|
              \le a^{-1/2}\bigl(3(2a\delta+2)+4/\delta\bigr)\le40 .
\]
Thus the fixed number of reflections works for every axis length.

\paragraph{Sharing the workspace.}
Divide the axes into consecutive blocks of four to seven. Such a
division is possible for large \(k\), since the number of axes tends
to infinity. Each block contains \(O(\log(2+k))\) address bits and
has only \(\operatorname{poly}(k)\) coordinates. Process its axes
in turn. For the chosen axis, the other padded axes give \(b_*\)
fibers, a power of two at least eight. Divide them into eight equal
batches. A batch of \(v=b_*/8\) fibers needs fewer than \(7va\)
borrowed coordinates, since each of its \(a\)-point transforms uses
\(2r-a<7a\) more coordinates. The other seven batches supply at
least \(7va\) coordinates. Allocate them disjointly to the active
fibers and run the circuits in parallel. Only the first \(a\) entries
of each active fiber change. Every borrowed entry is restored before
the next batch, even when it holds an earlier transform output.
The eight batches and at most seven axes keep the block depth
\(O(\log(2+k))\).

\paragraph{Three common layer types.}
In a layer of pair operations, complete its matching with identity
pairs. A block permutation routes the pairs to the edges of a single
address bit. The radix-two factorizations and the \(L(z)\) steps
require only pairs acting by \(H,L(z),L(z)^*\) or \(I\) once their
permutations and phases are separated. We can realize all pairs
with the same three Walsh rounds and pair-dependent phases.

For \(L(z)\), first form \(Q=H L(z)\). Every entry of \(Q\) has
modulus at least
\[
       \bigl(\sqrt{1-|z|^2}-|z|\bigr)/\sqrt2
           \ge(\sqrt3-1)/(2\sqrt2).
\]
Let \(c=|Q_{00}|\), \(s=|Q_{01}|\), \(\gamma=Q_{00}/c\), and put
\[
 \begin{aligned}
 D_{\rm in} &=\operatorname{diag}(1,Q_{01}/(is\gamma)),\\
 D_{\rm mid}&=\operatorname{diag}(c+is,c-is),\\
 D_{\rm out}&=\operatorname{diag}(\gamma,Q_{10}/(is)).
 \end{aligned}
\]
These are phase diagonals. The matrix \(HD_{\rm mid}H\) has \(c\)
on its diagonal and \(is\) off its diagonal. The unitary
\(D_{\rm out}HD_{\rm mid}HD_{\rm in}\) therefore agrees with \(Q\)
on the first row and first column. Orthogonality determines the fourth
entry, giving \(L(z)=HD_{\rm out}HD_{\rm mid}HD_{\rm in}\).
For \(L(z)^*\) replace \(z\) by \(\overline z\). For the identity
use \(Q=H\) in the same formulas. For a pair requiring \(H\), three
Walsh rounds with identity phases give \(H^3=H\).
In particular all divisions used to prepare the phases have
denominators bounded away from zero.

Undo the pair-routing permutation. The block schedules now consist
entirely of permutations, phase diagonals and Walsh rounds. Cycle
through the three types, advancing each block whenever its next type
occurs. Blocks that are idle in a slot act by the identity. Each
cycle advances every unfinished block, so the total number of global
layers is \(O(\log(2+k))\). Lemma~\ref{lem:controlled-blocks}
implements the permutations with empty external controls. A global
phase is the product of the local phases, and a global Walsh round
acts on at most one bit per block.

\paragraph{Streaming the coefficients.}
Every block has polynomially many coordinates. Its assignments of
borrowed positions, pair routings, radix-two permutations and layer
schedules can all be constructed in \(\operatorname{poly}(k)\) steps.
The same holds for its phase tables with \(B_0=\Theta(k)\) fractional
bits. The chirp and radix-two phases are roots of unity of
polynomial-size orders. The formulas for \(\lambda_\xi,D_2(\xi)\)
sum polynomially many such terms and scale by \(1/\sqrt a\).
Use range reduction and convergent sine and cosine series, for instance
with a Machin formula for \(\pi\), and bisection for positive square
roots. The arguments \(1-|z|^2\) are at least \(3/4\), and the phase
ratios above have uniformly bounded inverses of their denominators.
Only a bounded number of such operations follows each sum.
\(B_0+O(\log(2+k))\) fractional bits therefore suffice to make
the local tables accurate to \(2^{-B_0}\), by deterministic
polynomial-time fixed-tape computations.

To obtain a diagonal in tape order, extend address prefixes one
block at a time, from most to least significant. Start with the single
coefficient \(1\). For each coefficient on the current prefix tape,
scan the next block table, multiply by each of its entries, round,
and append the products to a fresh tape. Rewind the table between
prefixes and alternate the two prefix tapes. The length at least
doubles per pass, so fewer than \(2N\) products are written in all.
Each table rewind costs at most the scan it repeats. Multiplying
\(O_C(k)\)-bit words takes \(O_C(k\log(2+k))\) tape steps, giving
\(O_C(Nk\log(2+k))\) steps for the diagonal and the same bound to
apply it to the records.

If a local phase and a rounded complex product each have absolute
error at most \(u\), the prefix error increases from \(E\) to at
most \((1+u)E+2u\), since every ideal prefix has modulus one. After
at most \(k\) passes the error is \(O(ku)\) for \(ku\le1\). Choosing
the extra \(O(\log(2+k))\) guard bits gives absolute error \(2^{-Ck}\)
in the full coefficients while preserving the \(O_C(k)\)-bit width.
\end{proof}

\subsection{Rounding the convolution to the exact product}

Write \(A,B\) for the two padded digit arrays in residue order.
The subspace supported on valid addresses is invariant under
\(\mathcal U\) and its adjoint. There \(\mathcal U\) is the normalized
Fourier transform of a group of order \(S\). Thus, with \(\odot\)
denoting coordinatewise multiplication, character orthogonality gives
\[
 A*B=\sqrt S\,\mathcal U^*
                  \bigl((\mathcal U A)\odot(\mathcal U B)\bigr).
\]
After inverse CRT, entry \(t\) at a valid address is \([z^t]X(z)Y(z)\).
Indeed, \(2h-2<S\), so the ordinary product has no coefficient that wraps
around the cyclic index. It suffices to compute the displayed vector
to absolute Euclidean error less than \(1/2\).

\paragraph{Precision for one transform.}
Store the real and imaginary parts of every entry on signed fixed-point
grids. Suppose a forward or adjoint transform receives represented data
of norm at most \(2^{12k}\). Implement each selected-bit Walsh layer
with Lemma~\ref{lem:tape-walsh}, using an input norm allowance of
\(2^{13k}\) and an absolute Euclidean error allowance of \(2^{-50k}\).
Signed integer and fractional widths \(O_\rho(k)\) suffice. The proof
of the lemma permits enlarging these widths to any larger fixed
multiples of \(k\).

For a phase layer, compute the coefficients to absolute error at
most \(2^{-70k}\), and round each complex product to absolute error at
most \(2^{-70k}\). A diagonal perturbation has operator norm equal to
its largest entry error. On inputs of norm at most \(2^{13k}\), the
layer error, including rounding at \(N\) entries, is at most
\[
       2^{-70k}\bigl(2^{13k}+\sqrt N\bigr)<2^{-50k}
\]
for sufficiently large \(k\). Permutations move entire records
exactly. For an adjoint transform use inverse permutations and
conjugate phase tables. Every exact layer is unitary, so existing
errors pass through it with unchanged norm. Induction over the
\(O(\log k)\) layers keeps the boundary norm below \(2^{13k}\) and
the total error relative to the exact transform of the represented
input below \(2^{-40k}\). Lemma~\ref{lem:tape-walsh} also bounds
the intermediate widths within each Walsh layer.

\paragraph{Exact recovery.}
A digit array has \(h\) nonzero entries at most, each less than \(2^b\).
Its norm is at most \(\sqrt h\,2^b\le 2^{4k}\) for all sufficiently
large \(n\), since \(b/k\to1\). Let \(x=\mathcal U A\),
\(y=\mathcal U B\) and let \(x',y'\) be the rounded transform outputs.
By \(\|v\odot w\|_2\le\|v\|_2\|w\|_2\),
\[
 \|x'\odot y'-x\odot y\|_2
 \le \|x'-x\|_2\|y'\|_2+\|x\|_2\|y'-y\|_2.
\]
Rounding the products contributes at most \(\sqrt N\,2^{-70k}\).
The stored vector therefore has error at most \(2^{-30k}\) and norm
at most \(2^{12k}\). On this input the adjoint transform introduces
at most \(2^{-40k}\) further error. Compute \(\sqrt S\le2^{k/2}\) to
absolute error \(2^{-70k}\) by an integer square root in
\(\operatorname{poly}(k)\) bit steps. Even after scaling the existing
errors by \(\sqrt S\), multiplying by this rounded scalar and
rounding each entry, the final Euclidean error is less than \(1/2\).
All operands and temporary exact products fit in \(O_\rho(k)\) bits.
Rounding the real part of each valid entry to the nearest integer
now returns every convolution coefficient exactly. The bounded
initial range of input lengths again admits direct multiplication.

After inverse CRT, scan the valid addresses in increasing mixed-radix
rank. Each coefficient is less than \(h2^{2b}\). Add the incoming
carry, emit the low \(b\) bits, and pass the remaining high bits
to the next coefficient. This recurrence keeps the carry within
\(O(k)\) bits. Reverse the digit stream and trim or pad leading zeros
to write exactly \(2n\) bits in most-significant-first order. These
scans cost \(O(Nk)\).

\paragraph{Time and visited cells.}
Parameter generation, trial division, and construction of the local
permutation and phase tables take \(\operatorname{poly}_\rho(k)\)
steps. The full coefficient streams are generated as in
Lemma~\ref{lem:synchronized-short}. Each of the three transforms has
\(O(\log k)\) layers. A record permutation or selected Walsh layer
fits Equation~\eqref{eq:integer-budget}. A phase layer, including its
streamed coefficients, and the pointwise product each cost
\(O_\rho(Nk\log(2+k))+\operatorname{poly}_\rho(k)\) steps. Holding a
full stream while another input is processed costs \(O_\rho(Nk)\)
steps for each copy or rewind. The \(O(\log k)\) factor for the number
of layers is absorbed by the polylogarithmic factor in
Equation~\eqref{eq:integer-budget}. The number of tapes and the
alphabet depend only on \(\rho\).

A constant number of full streams occupy \(O_\rho(Nk)\) cells.
These include the two alternating prefix streams for the phase tables.
The permutations of Lemma~\ref{lem:controlled-blocks} sort at most
\(N\) exceptional records with \(O(k)\)-bit destination keys. Their
stream storage is \(O_\rho(Nk)\), and their interchanges have linear
space by Theorem~\ref{thm:interchange}. The Walsh layers visit
\(O_\rho(Nk)\) cells by Lemma~\ref{lem:tape-walsh}. The remaining
parameter and short-word computations visit
\(\operatorname{poly}_\rho(k)\) cells by their time bounds. All work
areas begin at marked positions and are reused after the heads return.
Hence the union of the visited cells on all tapes has size
\(O_\rho(Nk)+\operatorname{poly}_\rho(k)=O_\rho(n)\).
Finally \(Nk=\Theta(n)\) and \(k=\Theta(\log n)\). Our choice
\(\rho<1/2+\epsilon\) absorbs the remaining powers of \(\log k\)
in the time bound. This proves Theorem~\ref{thm:integer}.
